\documentclass[12pt]{article}
\usepackage[utf8]{inputenc}
\usepackage[T1]{fontenc}
\usepackage{hyperref}
\title{Tee Language Resources for K-12 Curriculum Development}
\author{Compiled by Claude Sonnet 5 (claude-sonnet-5) for Andr\'as Kornai}
\date{}
\begin{document}
\maketitle

\section{Introduction}
Tee (also Tai; ISO 639-3 tkq) is an East Ogoni language of Rivers State, Nigeria, in the Atlantic-Congo (Niger-Congo) family. It is only partly intelligible with Khana and its speakers regard it as a distinct language. Wikipedia gives about 310,000 speakers (2006). ASJP gives 485,000, so figures differ. Search found only general reference pages.

\section{Primary Educational Resources}
No primers, textbooks or school programmes were found in the search results.

\section{Linguistic and Reference Materials}
\begin{itemize}
\item Tee language, Wikipedia (sound system, seven oral and five nasal vowels, three tones): \url{https://en.wikipedia.org/wiki/Tee_language}
\item Khana language, Wikipedia (closely related language): \url{https://en.wikipedia.org/wiki/Khana_language}
\item Ogoni languages, Wikipedia: \url{https://en.wikipedia.org/wiki/Ogoni_languages}
\item ASJP word-list entry for Tee: \url{https://asjp.clld.org/languages/TEE}
\end{itemize}
No dictionary, grammar or official orthography guide was found. The tone and vowel details above come from secondary sources and should be checked against a primary one.

\section{Religious and Community Resources}
None found.

\section{Digital and Audio Resources}
None found.

\section{Curriculum Development Context}
The language is thought to be used by all in the ethnic community and is not known to be taught in schools. Almost no dedicated materials turned up, so this dossier cannot support curriculum design on its own. Next steps are Ogoni-language materials (Khana and Gokana have more documentation), Bible translation agencies, and Nigerian university linguistics departments.
\end{document}